\Needspace{8\baselineskip}
\Article{271}
Nun ist klar, dass alle ternären Formen ebenso wie die binären in \textit{Klassen} eingeteilt werden können, indem man äquivalente Formen derselben Klasse und nicht äquivalente Formen verschiedenen Klassen zuordnet. Formen mit verschiedenen Determinanten gehören also gewiss zu verschiedenen Klassen, und damit gibt es unendlich viele Klassen ternärer Formen. Ternäre Formen mit derselben Determinante bilden dagegen einmal weniger, einmal mehr Klassen. Als herausragende Eigenschaft dieser Formen ist aber anzusehen, dass \textit{alle Formen mit derselben gegebenen Determinante stets eine e\,n\,d\,l\,i\,c\,h\,e Anzahl von Klassen bilden.} Bevor wir diesen sehr wichtigen Satz ausführlicher entwickeln, müssen wir den folgenden wesentlichen Unterschied zwischen ternären Formen erläutern.

Einige ternäre Formen sind so beschaffen, dass durch sie ohne Unterschied positive und negative Zahlen dargestellt werden können, \textit{zum Beispiel} die Form $xx+yy-zz$. Sie werden daher \textit{indefinite Formen} heißen. Durch andere dagegen können keine negativen Zahlen, sondern nur positive dargestellt werden (außer der Null, die man erhält, indem man jede Unbestimmte $=0$ setzt), wie durch $xx+yy+zz$; deshalb werden sie \textit{positive Formen} heißen. Schließlich können durch andere keine positiven Zahlen dargestellt werden, wie durch $-xx-yy-zz$, weshalb sie \textit{negative Formen} genannt werden. Die positiven und negativen Formen werden zusammen \textit{definite Formen} heißen. Hier folgen nun allgemeine Kriterien, mit denen sich diese Eigenschaften der Formen unterscheiden lassen.

Multipliziert man die ternäre Form
\[
f=axx+a'x'x'+a''x''x''+2bx'x''+2b'xx''+2b''xx'
\]
\SourceLeaf{311}
mit der Determinante $D$ mit $a$ und bezeichnet die Koeffizienten der zu $f$ adjungierten Form wie in Artikel 267 mit $A$, $A'$, $A''$, $B$, $B'$, $B''$, so erhält man
\[
(ax+b''x'+b'x'')^2-A''x'x'+2Bx'x''-A'x''x''=g
\]
Multipliziert man nochmals mit $A'$, ergibt sich
\[
A'(ax+b''x'+b'x'')^2-(A'x''-Bx')^2+aDx'x'=h
\]
Daraus folgt sofort: Sind sowohl $A'$ als auch $aD$ negative Zahlen, so sind alle Werte von $h$ negativ. Durch die Form $f$ können dann offenbar nur solche Zahlen dargestellt werden, deren Vorzeichen dem von $aA'$ entgegengesetzt ist, \textit{also} mit dem von $a$ übereinstimmt oder dem von $D$ entgegengesetzt ist. In diesem Fall ist $f$ also eine definite Form, und zwar positiv oder negativ, je nachdem, ob $a$ positiv oder negativ ist oder ob $D$ negativ oder positiv ist.

Sind dagegen entweder beide Zahlen $aD$, $A'$ positiv oder die eine positiv und die andere negativ (keine $=0$), so sieht man leicht, dass $h$ durch geeignete Wahl der Größen $x$, $x'$, $x''$ sowohl positive als auch negative Werte annehmen kann. In diesem Fall kann also $f$ sowohl Werte mit demselben Vorzeichen wie $aA'$ als auch solche mit entgegengesetztem Vorzeichen annehmen und ist daher eine indefinite Form.

Für den Fall, dass $A'=0$, aber nicht $a=0$, gilt
\[
g=(ax+b''x'+b'x'')^2-x'(A''x'-2Bx'')
\]
Gibt man $x'$ einen beliebigen Wert (der aber nicht $=0$ ist) und wählt $x''$ so, dass $\frac{A''x'}{2B}-x''$ dasselbe Vorzeichen wie $Bx'$ hat (was offenbar möglich ist, da $B$ nicht $=0$ sein kann; denn daraus würde $BB-A'A''=aD=0$ und somit auch $D=0$ folgen, welchen Fall wir ausschließen), so ist $x'(A''x'-2Bx'')$ eine positive Größe. Daraus ist leicht zu sehen, dass $x$ so gewählt werden kann, dass $g$ einen negativen Wert erhält. Offenbar lassen sich diese Werte auch so wählen, dass sie, wenn gewünscht, alle ganzzahlig sind. Schließlich kann man bei beliebigen Werten von $x'$, $x''$ offenbar $x$ so groß wählen, dass $g$ positiv wird. Daraus folgt, dass die Form $f$ in diesem Fall indefinit ist.
\SourceLeaf{312}
Ist schließlich \(a=0\), so gilt
\[
f=a'x'x'+2bx'x''+a''x''x''+2x(b''x'+b'x'')
\]
Wählt man also \(x'\), \(x''\) beliebig, jedoch so, dass \(b''x'+b'x''\) nicht \(=0\) ist (was offenbar möglich ist, sofern nicht \(b'\) und \(b''\) zugleich \(=0\) sind; dann wäre aber \(D=0\)), so sieht man ohne Mühe, dass \(x\) so gewählt werden kann, dass \(f\) sowohl positive als auch negative Werte annimmt. Also ist \(f\) auch in diesem Fall eine indefinite Form.

Ebenso, wie wir hier aus den Zahlen \(aD\), \(A'\) die Beschaffenheit der Form \(f\) beurteilt haben, können wir auch \(aD\) und \(A''\) verwenden: \(f\) ist definit, wenn sowohl \(aD\) als auch \(A''\) negativ sind, und in allen übrigen Fällen indefinit. Auf genau dieselbe Weise kann man zu diesem Zweck auch die Zahlen \(a'D\) und \(A\), oder \(a'D\) und \(A''\), oder \(a''D\) und \(A\), oder schließlich \(a''D\) und \(A'\) betrachten.

Aus allem folgt, dass bei einer definiten Form die sechs Zahlen \(A\), \(A'\), \(A''\), \(aD\), \(a'D\), \(a''D\) negativ sind. Bei einer positiven Form sind dabei \(a\), \(a'\), \(a''\) positiv und \(D\) negativ; bei einer negativen dagegen \(a\), \(a'\), \(a''\) negativ und \(D\) positiv. Also zerfallen alle ternären Formen mit einer gegebenen positiven Determinante in negative und indefinite, alle mit einer negativen Determinante dagegen in positive und indefinite. Positive Formen mit positiver Determinante oder negative Formen mit negativer Determinante gibt es überhaupt nicht. --- Aus denselben Überlegungen sieht man leicht, dass einer definiten Form stets eine definite, und zwar \emph{negative}, einer indefiniten dagegen eine indefinite Form adjungiert ist.

Da alle durch eine gegebene ternäre Form darstellbaren Zahlen offenbar auch durch jede zu ihr äquivalente Form dargestellt werden können, sind die ternären Formen derselben Klasse entweder alle indefinit oder alle positiv oder alle negativ. Deshalb dürfen wir diese Bezeichnungen der Formen auch auf ganze Klassen übertragen.
